% Generated draft: requires mathematical and source review.
\documentclass[11pt]{article}
\usepackage[margin=25mm]{geometry}
\usepackage{fontspec}
\setmainfont{Latin Modern Roman}
\usepackage{amsmath,amssymb,mathtools,unicode-math}
\setmathfont{Latin Modern Math}
\usepackage{enumitem,needspace,xurl,hyperref}
\setlength{\parindent}{0pt}
\setlength{\parskip}{5pt}
\setlength{\emergencystretch}{3em}
\newcommand{\N}{\mathbb{N}}
\newcommand{\Z}{\mathbb{Z}}
\newcommand{\Q}{\mathbb{Q}}
\newcommand{\R}{\mathbb{R}}
\newcommand{\C}{\mathbb{C}}
\newcommand{\F}{\mathbb{F}}
\newcommand{\A}{\mathbb{A}}
\newcommand{\PP}{\mathbb{P}}
\newcommand{\E}{\mathbb{E}}
\newcommand{\eps}{\varepsilon}
\newcommand{\dd}{\,\mathrm{d}}
\begin{document}
\section*{UnsolvedMath ID: 30003002}
\textbf{Source problem code:} catalogue entry
\subsection{Definitions and mathematical statement}
Let \(n\) be a positive integer. A full-rank lattice in \(\R^n\) is a set \(L=A\Z^n\), where \(A\) is an invertible real \(n\times n\) matrix. A convex body is a compact convex subset of \(\R^n\) with nonempty interior. Under the usual origin-symmetry convention, a symmetric convex body satisfies \(K=-K\); the supplied record does not explicitly specify the center of symmetry. Define
\[
B_2^n=\left\{x\in\R^n:\sum_{i=1}^n x_i^2\leq 1\right\},\qquad tK=\{tx:x\in K\}\quad(t>0).
\]
For a finite set \(S\), \(|S|\) denotes its cardinality.

The supplied \emph{Strong Reverse Minkowski Conjecture} asks, for every symmetric convex body \(K\subset\R^n\) and every full-rank lattice \(L\subset\R^n\), whether
\[
|K\cap L|\leq M\bigl(O(\log n)K,L\bigr),
\qquad
|B_2^n\cap L|\leq M\bigl(O(\sqrt{\log n})B_2^n,L\bigr).
\]
These are schematic source expressions, not fully defined mathematical assertions: the supplied complete record gives no definition of \(M\).

Under the usual uniform asymptotic interpretation, the intended quantifiers would be the existence of constants \(C_1,C_2>0\) and an integer \(n_0\geq2\), independent of the body and lattice, such that for every integer \(n\geq n_0\), every permitted \(K\), and every full-rank \(L\subset\R^n\),
\[
|K\cap L|\leq M\bigl(C_1(\log n)K,L\bigr),
\qquad
|B_2^n\cap L|\leq M\bigl(C_2\sqrt{\log n}\,B_2^n,L\bigr).
\]
This interpretation is conditional and requires source confirmation. In particular, the definition of \(M\), the intended uniformity of the constants, and the treatment of small dimensions remain unresolved. The source's wording includes \(n=1\), where \(\log n=0\); no replacement such as \(\log(2n)\) is supplied.
\subsection{Short English statement}
The conjecture compares the number of lattice points in a symmetric convex body with an unspecified quantity \(M\) evaluated on a logarithmic dilation; for the Euclidean unit ball, the proposed dilation has order \(\sqrt{\log n}\). A precise complete statement requires the missing definition of \(M\) and clarification of the asymptotic quantifiers and small dimensions.
\subsection{Sources}
\begin{itemize}
\item Supplied UnsolvedMath local registry record, permanent ID \(30003002\), exact source code \texttt{catalogue entry}, in the Oberwolfach Reports Open Problems collection. Title: \emph{Strong Reverse Minkowski Conjecture}. The record supplies proposed year 2015 but no named proposer.
\item Supplied report citation: \emph{Convex Geometry and its Applications (2016)}, Oberwolfach Reports, pp.~3179--3269, DOI \texttt{10.4171/owr/2015/56}. \url{https://doi.org/10.4171/owr/2015/56}. The background describes these reports as records of mathematical workshops and their open-problem discussions.
\item Supplied literature source, labeled ``Strong reverse Minkowski program'': \url{https://arxiv.org/abs/1606.06913}.
\item Supplied literature source, labeled ``Recent reverse Minkowski theorem'': \url{https://annals.math.princeton.edu/2024/199-1/p01}.
\item Supplied literature source, labeled ``Original Oberwolfach conjecture'': \url{https://doi.org/10.4171/owr/2015/56}.
\item Literature assessment dated exactly \texttt{2026-08-21}: the supplied curation labels the problem partially solved and states that the full strong reverse Minkowski conjecture remains open, while reverse Minkowski theorems and lattice substructure results establish substantial weaker and special-case bounds. This is supplied dated curation derived from the OWR record, not an independently verified literature assessment; its status and cited claims require independent verification.
\item The registry labels the statement \texttt{corrected\_verified}, explaining that notation and publication prose were repaired without changing its substance. This is a supplied editorial label, not independent mathematical verification of this draft.
\end{itemize}
\end{document}
